% ==============================================================================
% Projeto de Sistema - Nome do Aluno
% Capítulo 1 - Introdução
% ==============================================================================
\chapter{Introdução}
\label{sec-intro}
\vspace{-1cm}

Este documento apresenta o projeto (\textit{design}) do sistema \emph{\imprimirtitulo}. O sistema tem como propósito gerenciar o ciclo completo de intermediação e contratação de serviços de
passeio para cães, permitindo o cadastro e validação de perfis de tutores, animais e passeadores, a realização de buscas filtradas por geolocalização e especialidade, o agendamento de passeios avulsos ou recorrentes, o monitoramento em tempo real do trajeto via GPS, a comunicação direta entre as partes e a efetivação de pagamentos, avaliações e registros fotográficos das atividades realizadas. 


Além desta introdução, este documento está organizado da seguinte forma: 
	a Seção~\ref{sec-plataforma} apresenta a plataforma de software utilizada na implementação do sistema;
	a Seção~\ref{sec-rnfs} apresenta a especificação dos requisitos não funcionais (atributos de qualidade), definindo as táticas e o tratamento a serem dados aos atributos de qualidade considerados condutores da arquitetura; 
	a Seção~\ref{sec-arquitetura} apresenta a arquitetura de software; 
	a Seção~\ref{sec-logica-negocio} apresenta o projeto da lógica de negócio dos subsistemas, com as componentes de domínio do problema e de gerência de tarefas; por fim, 
	a Seção~\ref{sec-ihc} apresenta o projeto da interação humano-computador, com a organização da Camada de Apresentação, o protótipo de alta fidelidade e o projeto da Camada de Interface com o Usuário.

